\section{The original finite contraction and its discarded indices}
\label{cta:finite}

\subsection{Fixed coordinates and the physical factors}
Write \(\theta=\pi/N\), \(\eta=(-1)^{N-1}\), and let
\(\mathsf T_N(z)\) be the monic polynomial with
\(\mathsf T_N(t+t^{-1})=t^N+t^{-N}\). The pair algebra at a vertex
has relation \(P_\sigma=\mathsf T_N+2\cos(2\pi\sigma)\), with
\(\sigma=0,s\), or \(-s\) prescribed by its oriented pair type.
The common vector-space basis is
\begin{equation}\label{cta:basis}
 e_k(z)=S_k^{-2}\prod_{j<k}(c_j-z),\qquad
 S_k=\prod_{j=1}^k2\sin(j\theta),\qquad
 c_j=2\cos((2j+1)\theta),\quad 0\le k<N.
\end{equation}
An empty product is one. A canonical representative has degree
less than \(N\), obtained using its actual relation \(P_\sigma\).
The different relations are not identified as algebras.

The closing functional and paired edge factors are
\begin{equation}\label{cta:pairing}
 \phi_\sigma(f)=\eta N^2[z^d]\rem_{P_\sigma}f,\qquad
 \phi_\sigma(e_i e_j)=\mathbf1_{i+j=d},
\end{equation}
\begin{equation}\label{cta:edge-factor}
 W_k(s)=\prod_{j=1}^k
  \left(1-\frac{\sin^2(\pi s/N)}{\sin^2(\pi j/N)}\right),
 \qquad W_kW_{d-k}=H_N^2.
\end{equation}
Every physical turn contributes its prescribed scalar
\(\gamma_e(s)\in\{1,q_N^{s^2/2},q_N^{-s^2/2}\}\).
Lowering its index with \eqref{cta:pairing} identifies the equal
indices on the two ends of the edge and leaves
\(\gamma_e H_N/W_k\), not its reciprocal.

Let \(p_v\) be the actual crossing power, including the carrier
specified by the original oriented diagram. The exact identity
of \cite[Theorem~7.1]{Meridian} is
\begin{equation}\label{cta:exact-source}
 \begin{split}
 Z_N(s)&=\sum_{\boldsymbol k\in\{0,\ldots,d\}^{E}}
       \prod_v\phi_{\sigma_v}
           \left(p_v\prod_{e\ni v}e_{k_e}\right)
       \prod_e\gamma_e(s)\frac{H_N(s)}{W_{k_e}(s)},\\
 F_N(q_N^s)&=\vartheta_{d+s}^{-w(D)}H_N(s)Z_N(s),
 \qquad \vartheta_{d+s}=q_N^{(s^2-d^2)/2}.
 \end{split}
\end{equation}
Thus cups, caps, physical turns, negative crossings, the pivotal
closure and the one reduced cut are already fixed finite data.

On every disk \(R<1\), all factors \(H_N/W_k\) and their inverses
are bounded independently of \(N,k\). For example split the factors
in \eqref{cta:edge-factor} according to
\(\min(j,N-j)\le J\) or \(>J\). The latter logarithms are bounded
by a summable constant times \(\min(j,N-j)^{-2}\).
For fixed \(J\) the former factors have no zeros on the disk;
their endpoint limits are \(1-s^2/j^2\), also zero-free.
Compactness, including the finitely many small \(N\), gives the
bound and its inverse. The same argument applies to \(H_N\).

\subsection{Bounds that retain the first normal}
Let \(m=\lfloor d/2\rfloor\). Indices \(k\le m\) will be called
low, the others high. For a canonical polynomial use the complete
norm
\[
 \|f\|_{\mathcal J}=
 \max\{|f(c)|,\ |f'(c)|:\ c\text{ a double node of }P_0;
                 \ |f(-2)|\text{ if }N\text{ is odd}\}.
\]
The derivative at the odd simple node is absent. All first
derivatives at double nodes are retained.

\begin{lemma}\label{cta:norms}
For fixed \(R<1\), a fixed number of factors and fixed integer
crossing powers, quotient multiplication, closing, and those
crossing powers have polynomial cost in \(\|\cdot\|_{\mathcal J}\).
Uniformly in all indices,
\begin{equation}\label{cta:basis-bound}
 \|e_k\|_{\mathcal J}\le C N^A
 \begin{cases}
   \exp(Ng(k/N)),&0\le k\le m,\\
   1,&m<k<N .
 \end{cases}
\end{equation}
The constants in these statements do not depend on \(s\) in its disk.
\end{lemma}
\begin{proof}
We give the conditioning and endpoint details because ignoring
derivatives would make the high-index assertion insufficient.
The distinct complete-node polynomial is
\[
 Q_N(z)=
 \begin{cases}
  2T_{N/2}(z/2),&N\text{ even},\\
  (z+2)\{U_m(z/2)-U_{m-1}(z/2)\},&N=2m+1.
 \end{cases}
\]
Here \(T_j,U_j\) are the ordinary Chebyshev polynomials.
Its interval norm and derivatives are polynomial in \(N\).
At its roots, \(|Q_N'|\ge N/2\) in the even case and at least
\(N/2\) in the odd case: for an odd interior root \(z=2\cos t\),
\(|Q_N'(z)|=N/(2\sin(t/2))\), while \(|Q_N'(-2)|=N\).
The mean-value formula bounds each Lagrange cardinal by a
polynomial in \(N\).

For a general polynomial of degree \(<N\), interpolate its values
by \(I\) of degree \(<\deg Q_N\). Write \(f-I=Q_N h\).
The values of \(h\) at the double nodes are
\((f'(c)-I'(c))/Q_N'(c)\). These determine \(h\), whose degree
is less than their number. Use \(Q_N\) for this interpolation at
even levels and \(Q_N/(z+2)\) at odd levels; the same cardinal
bound applies. Consequently the interval norm, the complete
normal norm, and the Chebyshev coefficient sum are polynomially
equivalent. One may bound derivatives without any conditioning
assumption: a degree-\(n\) Chebyshev expansion has coefficients
at most twice its interval norm, and
\(\|T_j'(z/2)\|\le j^2/2\).

Write \(C_j(z)=(t^j+t^{-j})/2\) when \(z=t+t^{-1}\).
The product identity \(2C_iC_j=C_{i+j}+C_{|i-j|}\) makes
the coefficient sum submultiplicative. Modulo \(P_\sigma\),
\[
 C_N=-\cos(2\pi\sigma),\qquad
 \rem C_{N+j}=-2\cos(2\pi\sigma)C_j-C_{N-j}\quad(0<j<N).
\]
Iteration bounds reduction of degree at most \(hN\) by
\((2\cosh(2\pi R)+1)^h\), for fixed \(h\).
Finally \([z^d]C_d=1/2\). These facts prove the asserted
polynomial costs of multiplication and \(\phi_\sigma\).

The actual crossing polynomial and its inverse are holomorphic
through the collisions in \(|\sigma|<1\), with the normals of
\cite[Proposition~6.1]{Meridian}. On a parameter circle of radius
\(R'\in(R,1)\), \(R'\ne1/2\), Lagrange interpolation has
denominator \(\sin(2\pi\sigma)\) bounded away from zero.
Its nodes have angular imaginary part \(O_{R'}(N^{-1})\).
The divided differences of \(\mathsf T_N\) and their derivatives
on that ellipse are \(O_{R'}(N^2)\) and \(O_{R'}(N^4)\):
differentiate \(2\cos(Nt)\), with endpoint limits interpreted
by its polynomial expression. The fixed-power eigenvalues,
including negative powers, are bounded there. Thus the canonical
crossing and its first derivative have polynomial interval norm.
Holomorphy and the maximum principle extend the bound through
all collisions to the disk. The mixed orders use their actual
fixed reference crossing and the bounded scalars of the same
proposition, so they satisfy the same conclusion.

For low \(k\), all values \(e_k(c_j)\) at distinct nodes are
nonnegative and bounded by \(e_k(-2)\). Interpolation by \(Q_N\)
therefore gives \(\|e_k\|_{\mathcal J}\le N^A e_k(-2)\).
Uniform sum--integral comparison of the midpoint cosine product
and sine factorial gives
\[
 \log e_k(-2)=Ng(k/N)+O(\log N),\qquad 0\le k\le m.
\]
This remains uniform at both endpoints: removing the first and
last mesh interval leaves a smooth logarithm, while each removed
logarithmic interval contributes \(O(\log N)\).

For high \(k\), put \(l=N-k\le\lfloor N/2\rfloor\) and
\(B_j(z)=\prod_{a<j}(c_a-z)\).
Reflection \(c_{d-a}=c_a\) gives
\(B_kB_l=(-1)^N P_0\). Thus \(e_k\) vanishes at every node.
At a double node indexed by \(0\le j<\lfloor N/2\rfloor\),
its derivative is zero unless \(j<l\), and otherwise
\begin{equation}\label{cta:high-normal}
 |e_k'(c_j)|
 =\frac{S_{l-1}^{\,2}}
 {2\sin((2j+1)\theta)S_{l+j}S_{l-1-j}}.
\end{equation}
To obtain this formula, compare quadratic coefficients in
\(B_kB_l=(-1)^NP_0\), using
\(P_0''(c_j)/2=N^2/(4\sin^2((2j+1)\theta))\),
\(|B_l'(c_j)|=S_{l+j}S_{l-1-j}/(2\sin((2j+1)\theta))\),
and \(S_kS_{l-1}=N\).
The denominator ratio
\(D_j=S_{l+j}S_{l-1-j}/S_{l-1}^2\) starts at
\(2\sin(l\theta)\) and is nondecreasing, since
\[
 \frac{D_{j+1}}{D_j}
   =\frac{\sin((l+j+1)\theta)}{\sin((l-1-j)\theta)}\ge1.
\]
The two angles sum to at most \(\pi\).
Equation \eqref{cta:high-normal} is consequently at most
\(1/(4\sin((2j+1)\theta)\sin(l\theta))\le N^2/16\).
The odd simple-node value is zero. This proves \eqref{cta:basis-bound}.
\end{proof}

\subsection{Discard before forming an angular sum}
Retain \(k_e/N\in I=[1/5,3/10]\) at every edge.
Each edge index occurs in two vertex factors of
\eqref{cta:exact-source}. Lemma~\ref{cta:norms} therefore bounds
the absolute rate of an assignment by
\(2\sum_{\text{low }e}g(k_e/N)\).
Direct differentiation gives
\[
 g'(1/4)=0,\qquad
 g''(x)=-\frac{4\pi}{\sin2\pi x}\le-4\pi.
\]
Since \(g(0)=g(1/2)=0\), concavity also gives \(2G_0\ge\pi/4\).
A low index outside \(I\) loses at least \(\pi/100\) from \(F_0\);
a high index loses at least \(2G_0\).
The number of assignments is polynomial in \(N\) for the fixed tree.
We have proved the full-source estimate
\begin{equation}\label{cta:discard}
 |\text{all discarded assignments}|
       \le C_R N^A e^{N(F_0-\eta_0)},\qquad \eta_0=\pi/200.
\end{equation}
The exponent loss is independent of tree size and powers; the
prefactor and starting level are not. This argument estimates
the actual vertex traces, rather than replacing a small spectral
norm by a positive coefficient cone.

On the retained indices, the edge sums commute with the polynomial
traces. Remove only the already recorded edge scalars and put
\begin{equation}\label{cta:packet}
 \begin{split}
 K_N(t,u;s)&=\sum_{k/N\in I}e_k(-2\cos t)e_k(-2\cos u)
                         \frac{H_N(s)}{W_k(s)},\\
 B_N(\boldsymbol\tau,s)&=\prod_{vw\in E}K_N(\tau_v,\tau_w;s).
 \end{split}
\end{equation}
This is an exact polynomial in each \(-2\cos\tau_v\).
Its separate degree may be larger than \(N-1\); no such bound
is assumed in the contours below. The functional \(\phi_{\sigma_v}\)
reduces the whole polynomial in its proper relation.
